% GENERATED FILE. Edit canonical lessons, then run npm run generate:books.
% canonical-source-hash: fnv1a64:ea76fa2794ca2791
% canonical-lessons: IT-C06-numeri-1-5, IT-C06-numeri-6-10

\chapter{Numbers One to Ten}
\label{ch:it-numbers-one-to-ten}
\hlchaptermodality{6}

\begin{chapteropening}
\textbf{By the end of this chapter:} \emph{I can count from one to ten in Italian and hear the old numbers still sitting inside settembre and dicembre.}

Count uno to dieci, pair sette, otto, nove and dieci with settembre, ottobre, novembre and dicembre, and hear Latin's -ct- flattened into the doubled -tt- of otto.
\end{chapteropening}

\section[uno, due, tre, quattro, cinque]{uno, due, tre, quattro, cinque — counting to five}

\label{lesson:IT-C06-numeri-1-5}



\begin{quote}
Of all the Romance sisters, \textbf{Italian kept the numbers closest to Latin} — it stayed in Italy, next to Rome, and wore the words down the least. Beside each you'll see the Spanish and French cousins for contrast.
\end{quote}

\begin{sounds}
\begin{itemize}
\raggedright
  \item \texttt{double-tt} — \textbf{quattro}: hold the double \emph{t} audibly, \emph{KWAT-tro} (Italian really doubles its consonants).
  \item \texttt{c-before-i} — \textbf{cinque} = \emph{CHEEN-kweh}: \emph{c} before \emph{i/e} is a \textquotedblleft{}ch\textquotedblright{} sound, but the \emph{-que} stays a hard \emph{kw}.
\end{itemize}
\end{sounds}

\begin{cousinweb}
\noindent
\begin{tabularx}{\linewidth}{@{}>{\raggedright\arraybackslash}X>{\raggedright\arraybackslash}X>{\raggedright\arraybackslash}X>{\raggedright\arraybackslash}X@{}}
\toprule
\textbf{Italian} & \textbf{$\leftarrow$ Latin} & \textbf{Spanish / French} & \textbf{English cousins} \\
\midrule
\textbf{uno} (1) & \emph{ūnus} & \emph{uno} / \emph{un} & \textbf{unit, union, unique} (and \emph{a/an}, \emph{one}) \\
\textbf{due} (2) & \emph{duo} & \emph{dos} / \emph{deux} & \textbf{duo, dual, duet, double} \\
\textbf{tre} (3) & \emph{trēs} & \emph{tres} / \emph{trois} & \textbf{trio, triple, triangle} \\
\textbf{quattro} (4) & \emph{quattuor} & \emph{cuatro} / \emph{quatre} & \textbf{quart, quarter, quatrain} \\
\textbf{cinque} (5) & \emph{quīnque} & \emph{cinco} / \emph{cinq} & \textbf{quintet, quintessence} \\
\bottomrule
\end{tabularx}

Notice \textbf{cinque} — Italian kept Latin \emph{quīnque} almost whole (\emph{-que} and all), while Spanish softened it to \emph{cinco} and French clipped it to \emph{cinq}. And \textbf{quattro} is why a musical foursome is a \emph{quartet} and four singers make a \emph{quattro}-anything. Italian is the sister who changed her Latin the least.
\end{cousinweb}

\begin{grammarlens}[title={Grammar Lens: uno $\to$ un / uno / una}]
\textbf{uno} is also \textquotedblleft{}a/an,\textquotedblright{} and it flexes by gender and the following sound: \textbf{un} (before most masc. nouns), \textbf{uno} (before tricky masc. clusters like \emph{s+consonant} or \emph{z-}), \textbf{una} (fem.). \emph{un caffè} \textquotedblleft{}a/one coffee,\textquotedblright{} \emph{una birra} \textquotedblleft{}a/one beer.\textquotedblright{} The others (\emph{due, tre…}) don't change.
\end{grammarlens}

\subsection*{Your turn}
Say these aloud:

\begin{itemize}
\raggedright
  \item \textquotedblleft{}uno, due, tre, quattro, cinque\textquotedblright{} — count up
  \item each with its English cousin — \textquotedblleft{}due/duo, tre/trio, quattro/quart, cinque/quintet\textquotedblright{}
  \item \textquotedblleft{}un caffè, una birra\textquotedblright{} — \emph{uno} $\to$ un/una before a noun
\end{itemize}

\begin{tcolorbox}[breakable,colback=teal!4,colframe=teal!35!black,title={Before you move on}]
Give 1–5 in Italian. (\emph{Uno, due, tre, quattro, cinque}.) Which Italian number kept Latin \emph{quīnque} most intact, versus Spanish and French? (\emph{cinque}, vs \emph{cinco} / \emph{cinq}.) What does \emph{un/una} mean besides \textquotedblleft{}one\textquotedblright{}? (\textquotedblleft{}A/an.\textquotedblright{}) Next: \textbf{6–10}, and the months hidden inside them.
\end{tcolorbox}
\section[sei, sette, otto, nove, dieci]{sei, sette, otto, nove, dieci — and the calendar's secret}

\label{lesson:IT-C06-numeri-6-10}



\begin{quote}
The rest of the count hides a calendar fact: \textbf{settembre–dicembre are Latin for 7, 8, 9, 10.} Italian even keeps the old Latin double consonants that show it.
\end{quote}

\begin{sounds}
\begin{itemize}
\raggedright
  \item \texttt{double-tt} — \textbf{sette} = \emph{SET-teh}, \textbf{otto} = \emph{OT-toh}: lean on the doubled \emph{t}, a giveaway that Latin's \emph{-pt-} and \emph{-ct-} both collapsed into \emph{-tt-}.
  \item \texttt{c-before-i} — \textbf{dieci} = \emph{DYE-chee} (the \emph{c} before \emph{i} $\to$ \textquotedblleft{}ch\textquotedblright{}).
\end{itemize}
\end{sounds}

\begin{cousinweb}
\noindent
\begin{tabularx}{\linewidth}{@{}>{\raggedright\arraybackslash}X>{\raggedright\arraybackslash}X>{\raggedright\arraybackslash}X>{\raggedright\arraybackslash}X@{}}
\toprule
\textbf{Italian} & \textbf{$\leftarrow$ Latin} & \textbf{Spanish twin} & \textbf{English cousins} \\
\midrule
\textbf{sei} (6) & \emph{sex} & \emph{seis} & \textbf{sextet, sextant} \\
\textbf{sette} (7) & \emph{septem} & \emph{siete} & \textbf{settembre} (the \emph{7th} month!), septet \\
\textbf{otto} (8) & \emph{octō} & \emph{ocho} & \textbf{ottobre} (8th), \textbf{octave, octopus} \\
\textbf{nove} (9) & \emph{novem} & \emph{nueve} & \textbf{novembre} (9th), novena \\
\textbf{dieci} (10) & \emph{decem} & \emph{diez} & \textbf{dicembre} (10th), \textbf{decimal, decade} \\
\bottomrule
\end{tabularx}

Watch the doubled consonants do historical work: Latin \emph{septem} had a \emph{-pt-} that Italian smoothed to \emph{-tt-} $\to$ \textbf{sette}; Latin \emph{octō} had a \emph{-ct-} that likewise became \emph{-tt-} $\to$ \textbf{otto}. Italian didn't drop those clusters (as French did in \emph{huit}) — it \textbf{assimilated} them into twins. So \emph{otto} and \emph{ottobre} still show the 8 side by side.
\end{cousinweb}

\begin{grammarlens}[title={Grammar Lens: why the months are \textquotedblleft{}off by two\textquotedblright{}}]
\emph{Settembre} means \textquotedblleft{}7th,\textquotedblright{} but it's the \textbf{9th} month. The old \textbf{Roman year began in March}, so \emph{September} really was the seventh; later \emph{luglio} and \emph{agosto} (July and August, for Julius and Augustus) pushed the counting months down two slots. The Latin numbers stuck to the labels. Every \emph{dicembre}, you say Roman \textquotedblleft{}\textbf{ten}\textquotedblright{} — just as in Spanish \emph{diciembre} and French \emph{décembre}.
\end{grammarlens}

\subsection*{Your turn}
Say these aloud:

\begin{itemize}
\raggedright
  \item \textquotedblleft{}sei, sette, otto, nove, dieci\textquotedblright{} — finish the count
  \item \textquotedblleft{}sette/settembre, otto/ottobre, nove/novembre, dieci/dicembre\textquotedblright{}
  \item \textquotedblleft{}octō $\to$ otto\textquotedblright{} — the \emph{-ct-} becoming a doubled \emph{-tt-}
\end{itemize}

\begin{tcolorbox}[breakable,colback=teal!4,colframe=teal!35!black,title={Before you move on}]
Give 6–10 in Italian. (\emph{Sei, sette, otto, nove, dieci}.) What happened to Latin's \emph{-ct-} in \emph{octō $\to$ otto}? (It assimilated to a doubled \emph{-tt-}.) Why is \emph{dicembre} Latin for \textquotedblleft{}ten\textquotedblright{}? (The Roman year began in March.) Next chapter builds on this toward time and days.
\end{tcolorbox}
